\section{What retraining blocks 0, 3 and 10 change}
\label{app:restore}
For blocks 0, 3 and 10 we retrained the full fine-tuned and the rank-64 shortcut models of each seed (the settings of Section~\ref{sec:fix};
all 18 pass the mitigation criterion) and patched between each pair right after the retrained block $b$, at layer $b+1$, on
the 2,000 \emph{opposite} reviews. The \emph{restore} puts the shortcut model's states into the retrained model, the
\emph{reverse} puts the retrained model's states into the shortcut model; both are shares of the same difference in the
name-aligned $d$ between the two models (Table~\ref{tab:restore}). After retraining block~0, the two names' states at
the name tokens differ \fourBAbRatioBZeroFullFTRange{} times as much as before (rank~64:
\fourBAbRatioBZeroLoraRSixtyFourRange).

\begin{table*}[!t]
\centering
\caption{Restore and reverse right after the retrained block (seed 0/1/2). The registered rule is $\ge 0.5$ in all three
seeds, for the name tokens and for the block-10 restore at the last token; the other last-token cells are
descriptive.}
\label{tab:restore}
\footnotesize
\begin{tabular}{llllll}
\toprule
Model & Block & Restore, name & Reverse, name & Restore, last & Reverse, last \\
\midrule
Full FT & 0 & \fourBRestoreNameBZeroFullFT & \fourBReverseNameBZeroFullFT & \fourBRestoreLastBZeroFullFT & \fourBReverseLastBZeroFullFT \\
Full FT & 3 & \fourBRestoreNameBThreeFullFT & \fourBReverseNameBThreeFullFT & \fourBRestoreLastBThreeFullFT & \fourBReverseLastBThreeFullFT \\
Full FT & 10 & \fourBRestoreNameBTenFullFT & \fourBReverseNameBTenFullFT & \fourBRestoreLastBTenFullFT & \fourBReverseLastBTenFullFT \\
Rank 64 & 0 & \fourBRestoreNameBZeroLoraRSixtyFour & \fourBReverseNameBZeroLoraRSixtyFour & \fourBRestoreLastBZeroLoraRSixtyFour & \fourBReverseLastBZeroLoraRSixtyFour \\
Rank 64 & 3 & \fourBRestoreNameBThreeLoraRSixtyFour & \fourBReverseNameBThreeLoraRSixtyFour & \fourBRestoreLastBThreeLoraRSixtyFour & \fourBReverseLastBThreeLoraRSixtyFour \\
Rank 64 & 10 & \fourBRestoreNameBTenLoraRSixtyFour & \fourBReverseNameBTenLoraRSixtyFour & \fourBRestoreLastBTenLoraRSixtyFour & \fourBReverseLastBTenLoraRSixtyFour \\
\bottomrule
\end{tabular}
\end{table*}

\textbf{Block 3.} Under rank~64 the restore at the name tokens holds in every seed, so the name tokens carry most of
block~3's fix; the reverse holds in two seeds of three (seed~1: \fourBReverseNameBThreeLoraRSixtyFoursOne). Under full
fine-tuning both are mixed. Restoring every position except the last brings back \fourBRestoreNotlastBThreeFullFTRange{}
(rank~64: \fourBRestoreNotlastBThreeLoraRSixtyFourRange), so block~3's fix sits on the review's tokens before the
answer, but not only on the name tokens.

\textbf{Block 10.} Under full fine-tuning the restore at the name tokens stays below 0.5 in every seed and the restore at
the last token is mixed; under rank~64 both are mixed. Neither position alone carries the fix, which fits block~10
being where the name's effect moves to the answer position (Section~\ref{sec:find}).

\section{Repair speed by block}
\label{app:budget}
Table~\ref{tab:speed} gives the per-seed values behind the speed pattern of Section~\ref{sec:how}.

\begin{table}[!htbp]
\centering
\caption{Fraction of the gap removed under limited training (seed 0/1/2). Groups are the mean over their blocks within
each seed. Full FT unless noted.}
\label{tab:speed}
\footnotesize
\setlength{\tabcolsep}{3pt}
\begin{tabular}{@{}llc@{}}
\toprule
Repair & Blocks & Removed \\
\midrule
Full-block, 100 steps & 0 & \budgetBZeroStepHundredFullFT \\
 & 2 & \budgetBTwoStepHundredFullFT \\
 & 4 & \budgetBFourStepHundredFullFT \\
 & 9 & \budgetBNineStepHundredFullFT \\
 & 10 & \budgetBTenStepHundredFullFT \\
Full-block, 100 steps, rank-64 model & 0 & \budgetBZeroStepHundredLoraRSixtyFour \\
 & 2 & \budgetBTwoStepHundredLoraRSixtyFour \\
 & 4 & \budgetBFourStepHundredLoraRSixtyFour \\
Weak fix, step $B$ & 0 & \uWeakStepBBZeroFullFT \\
 & 1--5 & \uWeakStepBBOneToFiveFullFT \\
 & 6--10 & \uWeakStepBBSixToTenFullFT \\
Weak fix, step $2B$ & 0 & \uWeakBZeroFullFT \\
 & 1--5 & \uWeakBOneToFiveFullFT \\
 & 6--10 & \uWeakBSixToTenFullFT \\
Weak fix, $2B$, rank-4 model & 0 & \uWeakBZeroLoraRFour \\
 & 1--5 & \uWeakBOneToFiveLoraRFour \\
Full-block, one epoch & 0 & \uRetrainBZeroFullFT \\
 & 1--5 & \uRetrainBOneToFiveFullFT \\
\bottomrule
\end{tabular}
\end{table}

\textbf{Early against late blocks.} The registered budget test retrained single blocks 0, 2, 4, 6, 8, 9, 10 and 11
with the settings of Section~\ref{sec:fix}, stopped at step 200, and predicted that at step 100 the mean fraction of
the gap removed by single blocks 9--11 exceeds that of single blocks 0, 2 and 4 by at least 0.20 in every seed. The
prediction failed under both full fine-tuning and rank~64 (per-seed differences in Tables~\ref{tab:prereg}
and~\ref{tab:prereg-rest}). Block~0 is the main reason (post hoc): it counts as early but is fast under both
(Table~\ref{tab:speed}). The weak-fix test registered two caveats: late blocks may be easier to fix in general, and
under LoRA the localized block comes from swaps whose own control failed.

\section{Why the LoRA swap control fails}
\label{app:lora-grid}
Two rank-64 LoRA models of different seeds differ in their adapter initialization and in their training path (data
order and dropout). We trained rank-64 models on a $3\times3$ grid of initialization seeds and data seeds (one of the
nine collapsed during training and was excluded by a rule fixed before any swap) and swapped blocks 9--11 between pairs
that differ only in the data seed (A), only in the initialization (B), or in both (C). Over all pairs, the swap fails to
keep the shortcut in \lFiveFailPairsCellA{} pairs of cell A and \lFiveFailPairsCellC{} of cell C, but in
\lFiveFailPairsCellB{} of cell B: the training path, not the initialization, makes late blocks incompatible across
runs, in line with the view that independently trained networks need not share a representation basis. The registered
test used three replicate pairs per cell (fraction kept: A \lFiveKeptCellARange, B \lFiveKeptCellBRange, C
\lFiveKeptCellCRange) and, with roughly one failing pair in three, read the original failure as a single-pair event.

\section{Where the adapters are placed}
\label{app:lora-where}
If LoRA adapters are placed only on blocks 0--5, the name's effect still reaches the answer position late, at layer
\lFourArrivalEarlyRange{} in every seed (first layer with $\Fl \ge 0.5$). The late takeover is therefore largely a
property of GPT-2's own later blocks, which are frozen here. With adapters only on blocks 6--11 it arrives at layer
\lFourArrivalLateRange, so the arrival also moves with where the shortcut is learned. The early-only shortcut is weaker
(gap \lFourGapEarlyRange, against \lFourGapLateRange{} for the late-only placement).

\section{A partly reliable name}
\label{app:partial}
To separate LoRA rank from shortcut strength, we tried to give rank~64 a shortcut as weak as rank~4's by making the name
agree with the label in only 90\%, 80\%, 70\% or 60\% of the named reviews. The seed-0 gaps were
\lSixGapPNinetysZero, \lSixGapPEightysZero, \lSixGapPSeventysZero{} and \lSixGapPSixtysZero; no level matched rank~4,
and the comparison was not run. The 90\% level fell just under the registered floor of 0.20, but the main reason is
the step size: rank~4's gaps (\gapLoraRFourRange) lie between those of the 100\% and 90\% levels, and the pilot tested
no level in between. Descriptively, at 90\% reliability only 50 contradicting reviews out of 10{,}000
reduce the gap from \gapLoraRSixtyFoursZero{} to \lSixGapPNinetysZero.

\section{A word-choice cue}
\label{app:wordcue}
We also planted a less obvious cue built from ordinary words: in a subset of the training reviews that use ``film'' or
``movie'', every occurrence was set to ``film'' in positive and ``movie'' in negative reviews. GPT-2 did not learn it.
With 5\% of the training reviews cued the seed-0 gap was \gapWordcueFivesZero, and with 10\% it was
\gapWordcueTensZero{} (clean twin \gapWordcueTenCleansZero), while the task itself was learned (acc(\textit{none})
\accNoneWordcueTensZero), so the find and fix analyses were not run. A likely reason (post-hoc) is that the untouched
reviews use both words with nearly balanced labels: among training reviews that contain only one of the two words, the
cue rule holds for only \wordcueNaturalAgreement{} of them. With those natural uses replaced by a neutral word the cue is
learned (Section~\ref{sec:othercues}).

\section{Registered predictions and outcomes}
\label{app:prereg-rest}
\begin{table*}[!tp]
\centering
\caption{Registered predictions behind the claims in the main text and the appendices, with their outcomes (values are
seed 0/1/2 unless noted). The other registered rules are listed in Table~\ref{tab:prereg-rest}. Decision rules were fixed before each run; the text reports effect sizes.
``Removes'' in a prediction means passing the mitigation criterion (Section~\ref{sec:setup}). A rule that needs all three seeds and fails means the effect was not
shown in every seed, not that it is absent. Close calls (a seed within about one test-set standard error of the
threshold, on the side that decides the outcome; post hoc) are noted in the Outcome column.}
\label{tab:prereg}
\setlength{\tabcolsep}{3pt}
\footnotesize
\begin{tabular}{p{4.5cm}p{3.5cm}p{5.9cm}p{3.0cm}}
\toprule
Prediction & Decision rule & Observed & Outcome \\
\midrule
Lens: true sentiment readable before a late takeover & lens reading rule & cal.\ acc(\textit{opposite}) at layer 9: \lensCalOppLNineFullFT & overthinking (late takeover) \\
Name reaches the answer only in blocks 9--11 & patching, $\Fl$ by layer & $\Fl(9) =$ \patchFLastLNineFullFT & holds (full FT): the name stays at its tokens until blocks 9--11 \\
Clean blocks 9--11 remove the shortcut (swap a) & mitigation criterion & \swapARemovedLateFullFT & holds (also blocks 10--11; no single block) \\
Swap control keeps the shortcut & kept $\ge 0.80$ & \swapCiiKeptLateFullFT & holds (seed 1: \swapCiiKeptLateFullFTsOne) \\
Shortcut blocks 9--11 create it in the clean twin (swap b) & created gap $\ge 0.20$ & FT \swapBCreatedLateFullFT; r64 \swapBCreatedLateLoraRSixtyFour & fails for FT (all 12 blocks create it); holds for r64 \\
Retraining any single block removes it & mitigation criterion & FT \retrainPassCountFullFT, r4 \retrainPassCountLoraRFour, r64 \retrainPassCountLoraRSixtyFour{} of 12 blocks & holds for FT and r64; not for r4 \\
Late blocks retrain better at 100 steps & late $-$ early $\ge 0.20$ & \budgetLateMinusEarlyFullFT & fails (seed 1: \budgetLateMinusEarlyFullFTsOne; holds in seed 0) \\
Swap and retrain rankings agree & Spearman $\ge 0.6$ & FT \spearmanSwapRetrainFullFT; r4 \spearmanSwapRetrainLoraRFour; r64 \spearmanSwapRetrainLoraRSixtyFour & FT, r4 hold; r64 fails \\
Counter-examples do the fixing & $\ge 3$ blocks pass with names, fail without & FT \ccCountedBlocksFullFT, r4 \ccCountedBlocksLoraRFour, r64 \ccCountedBlocksLoraRSixtyFour{} blocks & holds for all three \\
Retraining block 0 cuts the trigger: the fix sits at the name tokens (App.~\ref{app:restore}) & restore and reverse at the name tokens $\ge 0.5$ & restore FT \fourBRestoreNameBZeroFullFT, r64 \fourBRestoreNameBZeroLoraRSixtyFour; reverse FT \fourBReverseNameBZeroFullFT, r64 \fourBReverseNameBZeroLoraRSixtyFour & holds for both \\
Retraining block 3 or 10 acts on the name tokens (App.~\ref{app:restore}) & restore at the name tokens $\ge 0.5$ & block 3: FT \fourBRestoreNameBThreeFullFT, r64 \fourBRestoreNameBThreeLoraRSixtyFour; block 10: FT \fourBRestoreNameBTenFullFT, r64 \fourBRestoreNameBTenLoraRSixtyFour & r64 block 3 holds; FT block 10 fails; others mixed \\
Weak fix: the localized block is among the top 3 & rank $\le 3$ of 12 at step $B$ & rank FT \lThreeRankLocalizedFullFT, r4 \lThreeRankLocalizedLoraRFour, r64 \lThreeRankLocalizedLoraRSixtyFour & fails for all three \\
Weak fix: the best block is a top-3 swap block & best block in the swap top 3 & best FT \lThreeBestBlockFullFT, r4 \lThreeBestBlockLoraRFour, r64 \lThreeBestBlockLoraRSixtyFour & fails for all three \\
Weak fix follows the swap map & Spearman $\ge 0.6$ & FT \lThreeSpearmanFullFT; r4 \lThreeSpearmanLoraRFour; r64 \lThreeSpearmanLoraRSixtyFour & holds for FT only \\
Weak fix: block 0 is among the top 3 & rank $\le 3$ of 12 at step $B$ & rank FT \lThreeRankBZeroFullFT, r4 \lThreeRankBZeroLoraRFour, r64 \lThreeRankBZeroLoraRSixtyFour & fails for all three \\
Speed on two new FT models (test registered after seeds 0--2 were seen, before the new runs): block 0 beats blocks 1--5 & weak fix at step $B$; difference $\ge 0.10$ in both seeds & seed 3/4: \lSevenRuleOneFullFTsThree/\lSevenRuleOneFullFTsFour & holds \\
Speed on two new FT models: blocks 6--10 beat blocks 1--5 & mean difference $\ge 0.10$ in both seeds & seed 3/4: \lSevenRuleTwoFullFTsThree/\lSevenRuleTwoFullFTsFour & holds \\
Read-out at layer 9 removes the shortcut & mitigation criterion & \pSixRemovedLNineFullFT & holds \\
Random vector (same $\alpha$) removes $<$ half as much & ratio $< 0.5$ & ratio \pSixRandomRatioFullFT & fails (seed 1: \pSixRandomRatioFullFTsOne; holds in seed 2) \\
Equal dose: layer 9 best & $l^*=9$ removes most & \pSixBEqLNineRemovedFullFT & holds (seed 1: accuracy drop \pSixBEqLNineDropCountFullFTsOne{} of 40 allowed reviews) \\
Random vector (own $\alpha$) removes $<$ half as much & ratio $< 0.5$ & ratio \pSixBRandomRatioFullFT & holds \\
Cue-free $\alpha$ still removes the shortcut & mitigation criterion & \pSixBCueFreeRemovedFullFT & holds (seed 1: accuracy drop \pSixBCueFreeDropCountFullFTsOne{} of 40 allowed reviews) \\
LoRA swap control keeps the shortcut & kept $\ge 0.80$ & r4 \swapCiiKeptLateLoraRFour; r64 \swapCiiKeptLateLoraRSixtyFour & fails (r4 seed 1: \swapCiiKeptLateLoraRFoursOne; r64 only in seed 0) \\
LoRA control failure: init or training path (App.~\ref{app:lora-grid}) & 3 replicate pairs per cell & A \lFiveKeptCellA; B \lFiveKeptCellB; C \lFiveKeptCellC & single-pair event (C pair 1: \lFiveKeptCellCsZero) \\
LoRA on blocks 0--5: arrival layer (App.~\ref{app:lora-where}) & $\ge 7$: GPT-2's structure & \lFourArrivalEarly & GPT-2's structure \\
Partly reliable name matches rank 4 (App.~\ref{app:partial}) & a level with gap 0.20--0.45 & best \lSixGapPNinetysZero{} (seed 0) & no level; stopped (just under 0.20) \\
Word-choice cue is learned (App.~\ref{app:wordcue}) & gap $\ge 0.20$ & \gapWordcueFivesZero{} (5\%), \gapWordcueTensZero{} (10\%), seed 0 & not learned; stopped \\
Memorized labels: no single block removes them & no block passes & block \wlPassingBlockNoisy{} removes \wlRetrainSingleNoisy & fails (block \wlPassingBlockNoisy{} passes) \\
Pythia: lens reading (mean over seeds) & takeover layer & \pyTakeoverMeanPythia & decided early \\
Pythia: retraining any single block removes it & mitigation criterion & \pyRetrainSingleCountPythia{} of 12 blocks & holds \\
Pythia: swap and retrain rankings agree & Spearman $\ge 0.6$ & \pyRetrainSpearmanPythia & fails (seed 2: \pyRetrainSpearmanPythiasTwo) \\
\bottomrule
\end{tabular}
\end{table*}

Table~\ref{tab:prereg} lists the registered predictions behind the claims in the main text and the appendices;
Table~\ref{tab:prereg-rest} lists the other registered rules. Training checks, run
gates and selection rules (learning rate, step budget $B$, the partly-reliable level, Pythia's names) decide how an
experiment is run rather than what it finds, and are not listed.

\begin{table*}[!tp]
\centering
\caption{Registered rules not in Table~\ref{tab:prereg}, with their outcomes (values are seed 0/1/2 unless noted). Same
reading as Table~\ref{tab:prereg}: a rule that needs all three seeds and fails means the effect was not shown in every
seed, and ``removes'' means passing the mitigation criterion.}
\label{tab:prereg-rest}
\setlength{\tabcolsep}{3pt}
\footnotesize
\begin{tabular}{p{4.5cm}p{3.5cm}p{5.9cm}p{3.0cm}}
\toprule
Prediction & Decision rule & Observed & Outcome \\
\midrule
Clean blocks of another seed remove the shortcut (swap c-i) & mitigation criterion & same block sets as case (a) & holds at blocks 9--11 and 10--11 \\
Blocks 9--11 are late-specific (swap a) & removal $\ge 2\times$ best other 3-block window & FT ratio \twoBLateRatioFullFT & fails for FT; holds for r4 and r64 \\
Name embeddings carry part of the shortcut (swap b) & name rows add $\ge 0.10$ to blocks 9--11 & FT \twoBEmbNamesExtraFullFT & fails \\
Attention or MLP half carries it (swap a) & half removes $\ge 2/3$ of the whole & blocks 9--11: attention \twoBAttnShareFullFT, MLP \twoBMlpShareFullFT & both needed at 9--11; attention alone at block 11 \\
Retraining on shortcut data keeps it (control) & kept $\ge 0.80$ & blocks 9--11: FT \retrainCtrlKeptFullFT; r4 \retrainCtrlKeptLoraRFour; r64 \retrainCtrlKeptLoraRSixtyFour & holds for all three \\
Late blocks retrain better at 100 steps, r64 & late $-$ early $\ge 0.20$ & \budgetLateMinusEarlyLoraRSixtyFour & fails \\
LoRA lens: same reading as full FT & lens reading rule & cal.\ acc(\textit{opposite}) at layer 9: r4 \lensCalOppLNineLoraRFour; r64 \lensCalOppLNineLoraRSixtyFour & holds for both (overthinking) \\
LoRA swap map matches full FT & Spearman $\ge 0.6$ & r4 \spearmanSwapVsFullFTLoraRFour; r64 \spearmanSwapVsFullFTLoraRSixtyFour & r4 holds; r64 fails \\
Weak-fix control keeps the shortcut & kept $\ge 0.80$ & FT \lThreeCtrlKeptFullFT; r4 \lThreeCtrlKeptLoraRFour; r64 \lThreeCtrlKeptLoraRSixtyFour & holds for all three \\
Speed test on two new FT models: control keeps the shortcut & kept $\ge 0.80$ & seed 3/4: \lSevenCtrlKeptFullFTsThree/\lSevenCtrlKeptFullFTsFour & holds \\
LoRA on blocks 6--11 only learns the shortcut & gap $\ge 0.20$ & \lFourGapLate & holds \\
Which LoRA blocks, not how many, sets the shortcut's strength & gap differs by $\ge 0.10$ from both alternate sets & early \lFourGapEarly, late \lFourGapLate, even \lFourGapEven, odd \lFourGapOdd & mixed for both sets \\
Retraining block 10: the fix acts at the answer position & restore at the last token $\ge 0.5$ & FT \fourBRestoreLastBTenFullFT; r64 \fourBRestoreLastBTenLoraRSixtyFour & mixed for both \\
Retraining block 3 or 10: the name tokens alone carry the fix & reverse at the name tokens $\ge 0.5$ & block 3: FT \fourBReverseNameBThreeFullFT, r64 \fourBReverseNameBThreeLoraRSixtyFour; block 10: FT \fourBReverseNameBTenFullFT, r64 \fourBReverseNameBTenLoraRSixtyFour & mixed for all four \\
Partly reliable name: location matters at rank 64 (App.~\ref{app:partial}) & $\ge 3$ single blocks fail & -- & not run (no level qualified) \\
Word-choice cue: find and fix rules (App.~\ref{app:wordcue}) & as for the name & -- & not run (cue not learned) \\
Memorized labels: lens shows a late takeover & lens reading rule & takeover layer \wlTakeoverNoisy & holds \\
Memorized labels: swap control keeps them & kept $\ge 0.80$ & blocks \wlControlBlocksNoisy: \wlControlKeptNoisy & holds \\
Memorized labels: retraining on the noisy labels keeps them (control) & kept $\ge 0.80$ & blocks 9--11 and all blocks & holds \\
Memorized labels: swap and retrain rankings agree & Spearman $\ge 0.6$ & \wlSpearmanNoisy & fails \\
Pythia: clean blocks 9--11 remove the shortcut & mitigation criterion & \pySwapRemovedLatePythia & fails (block 0 alone removes \pySwapRemovedBZeroPythia) \\
Pythia: shortcut blocks 9--11 create it (swap b) & created gap $\ge 0.20$ & -- & fails \\
Pythia: swap control keeps the shortcut & kept $\ge 0.80$ & blocks 9--11: \pySwapCiiKeptLatePythia & holds \\
Pythia: 3-block windows (swap a, b, c-ii) & as for GPT-2 & -- & blocks 0--2 remove and create it; no window passes the control in every seed; name rows add nothing \\
Pythia: retraining block sets removes it & mitigation criterion & windows 0--2, 3--5, 6--8, 9--11, 10--11; all: \pyRetrainRemovedAllPythia & holds for all six \\
Pythia: retraining on shortcut data keeps it (control) & kept $\ge 0.80$ & blocks 9--11: \pyRetrainControlKeptLatePythia; all: \pyRetrainControlKeptAllPythia & holds for both \\
\bottomrule
\end{tabular}
\end{table*}
